\chapter{Conclusion}\label{ch:conclusion}
\graphicspath{{../Figures/80-conclusion/}{figs/80-conclusion/}}

This dissertation moved the biomimetic humanoid robot project from an unvalidated design concept to a quantitatively characterized actuator set and experimentally tested knee-torque models, and it defined the path from those experimental foundations to neuromechanical control.

Chapter~\ref{ch:muscles} produced equations for isometric force in Festo \qtylist{10;20}{\mm}-diameter BPAs. The resting-length dependence $F_{620}(l_0)$ corrects the manufacturer tool's opposing trend at \qty{620}{\kPa}, and the nondimensional surface $F^*(\epsilon^*,P^*)$ uses three diameter-dependent coefficients. Together, these fits support actuator selection and route evaluation across resting lengths before prototyping, within the tested isometric conditions.

At the joint level (Chapter~\ref{ch:torque}), the BPA model was combined with design factors for effective length, series compliance, and actuator bending. Identification on the pinned knee improved torque prediction, and the shared design factors also improved predictions on the biomimetic four-bar knee. The improved model supports route redesign: the optimized two-BPA flexor and cam-extensor configurations are predicted to meet their human-scaled torque targets across the design range. In the isometric tests, the redesigned flexor and cam-assisted extensor exceeded their respective human-muscle reference torques at the measured knee angles. Agreement with the pressure-specific predictions varied with knee angle.

The SADb application in Chapter~\ref{ch:database} makes physiological evidence searchable by sensory pathway, animal or human study context, and modeling relevance. Linked records, curation notes, and citation navigation connect findings to their sources. This resource helps researchers compare evidence and select pathways for neuromechanical models.

Chapter~\ref{ch:sns} brings together the Li--Deng/W2L and Gait2392 models, SNS-Simscape blocks, the MuJoCo--Simulink interface, and procedures for tuning and testing controllers. Matched circuit implementations reproduced the reference activity, and the interfaces carried motor commands and returned sensory measurements. In the simulated arm, stretch feedback recruited the biceps and reduced sag under load. The feet of the W2L robot model contacted the ground while a virtual harness supplied an upward force to its body. Under those conditions, disabling force feedback through the Ib pathways alone preserved alternating heel strikes. Disabling heel and toe contact, pooled extensor-force inputs, and their associated spinal connections together disrupted that timing while the model continued forward. Measuring both heel-strike timing and forward displacement under the same drive and support distinguished the effect on coordination from the effect on progression. General controller transfer and independent balance remain open.

These actuator and joint models, the sensory-afferent application, and the neuromechanical tools make design assumptions testable: force against measurement, routing against torque requirements, and sensory pathways against controlled changes in movement. The future studies in Chapter~\ref{ch:overall} extend these tests toward dynamic actuation, independent balance, and hardware. This work provides an experimental foundation for biomimetic robot design and for studying how neural control interacts with mechanical structure.
